\section{Limitaciones de la implementaci\'on}

La limitaci\'on mas importante de esta implementaci\'on es que no se puede 
agregar vectores al \'indice mas de una vez. Esto no es sumamente importante 
(al menos para los ejemplos vistos), pero es una mejora pendiente.

\subsection*{Diferencias con el art\'iculo de referencia\cite{pqnns}.}

El art\'iculo describiendo PQ para NNS\cite{pqnns} incluye varias mejoras al 
algoritmo b\'asico.

\begin{itemize}
  \item\textbf{Busqueda con MaxHeap:} 
    Para la busqueda de los $k$ vectores m\'as cercanos, el art\'iculo describe
    el uso de la estructura de datos MaxHeap para mayor velocidad de
    inserci\'on ordenada.
  \item\textbf{Calculo de distancias con residuales:} 
    Se usa un m\'etodo modificado para el c\'alculo de distancias, que requiere 
    el c\'alculo de un vector residual (\(r(y) = y - q_c(y)\)) para un menor uso 
    de memoria. No logr\'e encontrar esto en la implementaci\'on de FAISS 
    tampoco.
  \item\textbf{Tama\~no de los codigos cuantificados:} 
    Uno de los parametros de \texttt{index\_pq\_init} es la cantidad de bits 
    usados para el indexado a los centroides, que se usa para calcular la 
    cantidad de centroides calculados. Sin embargo, si tenemos \texttt{n\_bits = 
    8}, lo ideal ser\'ia guardar los codigos en 8 bits, pero por simplicidad los 
    guardamos siempre en \texttt{int}s (32 bits).
\end{itemize}
